\section{RKNN-FS}

\begin{frame}
\frametitle{Parâmetros}
\begin{itemize}
    \item \textit{r}: quantidade de KNNs usados
    \item \textit{q}: taxa de redução de variáveis no primeiro estágio
    \item \textit{d}: quantidade de variáveis eliminadas no segundo estágio
    \item \textit{m(p)}: quantidade de variáveis consideradas por cada KNN, tal que \textit{p} é a quantidade de variáveis restantes e $ m(p) < p $
\end{itemize}
\end{frame}

\begin{frame}
\frametitle{Algoritmo}
\begin{itemize}
    \item \textit{r} KNNs são criados, cada um com \textit{m(p\textsubscript{0})} variáveis, sendo \textit{p\textsubscript{0}} o total de variáveis
    \item A cada iteração,
    \begin{itemize}
        \item os KNNs são treinados sobre suas \textit{m(p)} variáveis
        \item A métrica \textit{feature support} é calculada para cada variável
        \item As variáveis são ordenadas conforme a métrica
        \item As \textit{p} variáveis mais relevantes são mantidas, com \textit{p} computado conforme o estágio
    \end{itemize}
    \item As variáveis restantes ao fim de todas as etapas de eliminação são passadas para o KNN de inferência
\end{itemize}
\end{frame}

\subsection{\textit{Feature Support}}

\begin{frame}

\begin{itemize}
    \item Cada variável \textit{f} é atribuída a \textit{M} classificadores, que formam o conjunto \textit{C(f)}. Cada classificador recebe o conjunto \textit{F(c)} de variáveis
    \item Quando um classificador \textit{c} acerta sua predição, as variáveis $f \in F(c)$ recebem uma pontuação positiva
    \item A pontuação de cada variável \textit{f} é a média das pontuações atribuídas por cada classificador $c \in C(f)$
    \item Variáveis mais relevantes para a classificação têm pontuação maior
\end{itemize}

\end{frame}

\subsection{Seleção de Variáveis}

\begin{frame}
\frametitle{Remoção Geométrica}
\begin{itemize}
    \item A cada iteração, a metade menos relevante das variáveis é descartada
    \item Após todas as iterações, a quantidade de variáveis que resulta na maior precisão é encontrada
    \item As variáveis presentes na iteração anterior a esta são passadas para a próxima etapa
\end{itemize}
\end{frame}

\begin{frame}
\frametitle{Remoção Linear}
\begin{itemize}
    \item A cada iteração, um número fixo de variáveis é removido
    \item Após todas as iterações, o conjunto de variáveis mais relevantes é selecionada para o modelo
\end{itemize}
\end{frame}
